\section{Clearing local factors by a reciprocal fourth power}
This proof checkpoint records an algebraic reduction whose complete
coefficient reconstruction is in progress. It asserts no polynomial
count until that reconstruction and the exceptional cases are checked.

For a nonzero rational weight $q$, a factor
$T=1+q^5C^2+q^2D^2+q^4B^4$ can instead be imposed as
\[
 r_4T=r_4+qC^2+r_2D^2+(B^2)^2,
 \qquad qr_1=1,\quad r_2=r_1^2,\quad r_4=r_2^2.
\]
The field squares $C^2,D^2,B^2$ are already stored in the source.
Thus each field fourth-power coordinate $B^4$ is removed while the
three scalar weights replace three existing scalar weights.
The five inherited comparison factors remove ten rational coordinates.
At a consumed prime with $v(q)=-\nu$ for odd $\nu>0$, the fifth-power
comparison term gives precision three. In the mixed factor containing
$I^4$ and an even-valuation reciprocal abscissa $A$, use $q^5I^4$
and $q^3A^2$; their valuations are respectively $-\nu$ and $\nu$
modulo four. This preserves separation of the local terms.

In the additional-abc construction, the three reciprocal weights
$r_2,r_3,r_4$ satisfy
\[
 qr_3=r_2,\qquad qr_4=r_3,\qquad r_2^2=r_4.
\]
The nonzero component gives $r_j=q^{-j}$. On the zero component,
the two cleared factor equations force both numerator reciprocal
abscissas to vanish. The existing fixed-difference chart then gives
both numerator points equal to the identity, contradicting their
nonzero difference. A full polynomial proof of this exclusion is
required in the final construction.
